\documentclass[11pt,a4paper]{article}
% Estilo comum dos PDFs individuais da Sprint 1.
\usepackage[utf8]{inputenc}
\usepackage[T1]{fontenc}
\usepackage[brazilian]{babel}
\usepackage[a4paper,margin=2.35cm]{geometry}
\usepackage{xcolor}
\usepackage{hyperref}
\usepackage{enumitem}
\usepackage{listings}
\usepackage{parskip}
\usepackage{needspace}
\definecolor{azul}{RGB}{20,63,112}
\hypersetup{colorlinks=true,linkcolor=azul,urlcolor=azul}
\setlist[itemize]{leftmargin=*,itemsep=3pt}
\lstset{basicstyle=\ttfamily\small,frame=single,breaklines=true,columns=fullflexible}
\sloppy
\newcommand{\codigo}[1]{\texttt{#1}}
\newcommand{\titulo}[2]{\begin{center}{\LARGE\bfseries #1\par}\vspace{4pt}{\large #2\par}\vspace{8pt}\rule{0.8\textwidth}{0.4pt}\par\vspace{5pt}{\small Matheus Moreli Durão Silva -- Matrícula 12621CBS217\par Cibersegurança - Bacharelado - Noturno\par}\end{center}}
\newcommand{\cardfonte}[1]{\par\smallskip\noindent\textit{Referência: }\href{#1}{Trello - enunciado e checklist do card}.}

\begin{document}\titulo{S1\_03 - Critérios de aceitação}{Como saber se uma funcionalidade está pronta}
\section*{O que são critérios de aceitação}
Um requisito diz o que precisa existir; o critério de aceitação diz como verificar se foi atendido. Ele deve ser objetivo, observável e testável. Assim, ``o cadastro deve funcionar'' é vago; ``ao informar ID, nome, responsável, localização e tipo válidos, o ativo deve aparecer na listagem'' é um critério verificável.
\section*{Critérios definidos}
\textbf{Cadastro de ativo:} campos obrigatórios não podem ficar vazios; o ID deve ser numérico; um ID já utilizado deve gerar mensagem de impedimento; após confirmar, a listagem deve mostrar os dados. \textbf{Registro de vulnerabilidade:} deve ser escolhida uma severidade válida; a vulnerabilidade deve ser associada a um ativo existente; a consulta posterior deve exibir descrição e severidade.
\section*{Casos de teste manual}
\begin{itemize}
\item Caso válido: cadastrar ID 10, nome ``Notebook do laboratório'', responsável ``Equipe TI'' e localização ``Sala 3''. Esperado: confirmação e item visível ao listar.
\item Caso inválido: tentar cadastrar o mesmo ID 10 ou deixar o nome vazio. Esperado: aviso de validação e nenhum novo ativo criado.
\item Vulnerabilidade: associar ``senha padrão ativa'' ao ID 10 com severidade Alta. Esperado: registro visível ao consultar vulnerabilidades.
\end{itemize}
\section*{Relação com o projeto}
As validações e mensagens do \codigo{app.py} não existem por enfeite: elas satisfazem critérios. Em uma entrega, testar significa executar um caso com entrada conhecida, comparar a saída ao resultado esperado e registrar se passou ou falhou. A definição de pronto (\textit{Definition of Done}) desta etapa é: requisito implementado, entradas inválidas tratadas e resultado conferido.
\section*{Checklist demonstrado}
Conceito de critérios, relação com requisitos, validação, definição de pronto e exemplos aplicados ao Trabalho 1 foram contemplados neste recorte.
\cardfonte{https://trello.com/c/MRXQneEZ}
\end{document}
